\section{مدل‌سازی و تعریف مسئله}

\subsection{مدل برنامه‌ریزی خطی عدد صحیح (\lr{ILP})}\label{subsec:lr_ilp}
فرض کنید $I=\{1,\ldots,n\}$ مجموعه اقلام، $w_i>0$ وزن قلم $i$، و $C>0$ ظرفیت هر سطل باشد. گراف بدون جهت $G=(I,E)$ روابط چسبندگی را نمایش می‌دهد؛ یال $\{i,k\}\in E$ یعنی دو قلم باید در یک سطل قرار گیرند. چون هر جواب موجه حداکثر از $n$ سطل استفاده می‌کند، مجموعه سطل‌های نامزد را $J=\{1,\ldots,n\}$ می‌گیریم. متغیر $x_{ij}$ نشان می‌دهد قلم $i$ در سطل $j$ قرار دارد و $y_j$ استفاده‌شدن سطل $j$ را نشان می‌دهد:
\begin{align}
    \min \quad & \sum_{j\in J} y_j \label{eq:obj} \\
    \text{s.t.}\quad
    & \sum_{j\in J} x_{ij}=1 && \forall i\in I, \label{eq:assign} \\
    & \sum_{i\in I} w_i x_{ij}\leq C y_j && \forall j\in J, \label{eq:capacity} \\
    & x_{ij}=x_{kj} && \forall \{i,k\}\in E,\ \forall j\in J, \label{eq:cohesion} \\
    & x_{ij}\in\{0,1\} && \forall i\in I,\ \forall j\in J, \label{eq:x_binary} \\
    & y_j\in\{0,1\} && \forall j\in J. \label{eq:y_binary}
\end{align}
قید \eqref{eq:assign} هر قلم را دقیقاً یک‌بار تخصیص می‌دهد. قید \eqref{eq:capacity} هم ظرفیت را اعمال می‌کند و هم، به علت مثبت‌بودن وزن‌ها، اجازه نمی‌دهد قلمی در سطل با $y_j=0$ قرار گیرد. برای هر یال، قیود \eqref{eq:cohesion} بردار تخصیص دو رأس را برابر می‌کنند؛ در نتیجه برابری در امتداد هر مسیر منتشر می‌شود و همه رأس‌های یک مؤلفه همبند دقیقاً در یک سطل قرار می‌گیرند. بنابراین لازم نیست برای تمام جفت‌رأس‌های هر مؤلفه قید جداگانه افزوده شود.

\subsection{انقباض مؤلفه‌ها و هم‌ارزی دقیق با \lr{BPP}}\label{subsec:lr_component_based}
مؤلفه‌های همبند $G$ را با $\mathcal{C}$ نشان می‌دهیم و برای هر $c\in\mathcal{C}$ یک اَبَرقلم با وزن
\[
    W_c=\sum_{i\in c} w_i
\]
می‌سازیم. این انقباض فقط یک تقریب یا ابتکار نیست، بلکه تبدیل دقیق مسئله است:
\begin{itemize}
    \item در هر بسته‌بندی موجه مسئله اصلی، قیود \lr{must-link} تمام اقلام $c$ را در یک سطل قرار می‌دهند؛ پس جایگزینی آن‌ها با اَبَرقلمی به وزن $W_c$ ظرفیت و تعداد سطل‌ها را تغییر نمی‌دهد.
    \item برعکس، از هر بسته‌بندی موجه اَبَراقلام می‌توان با قرار دادن همه اقلام $i\in c$ در سطل اَبَرقلم $c$، یک بسته‌بندی موجه مسئله اصلی با همان تعداد سطل ساخت.
\end{itemize}
پس مجموعه جواب‌ها، پس از نادیده‌گرفتن تمایز داخلی هر مؤلفه، متناظر است و مقدار بهینه دقیقاً حفظ می‌شود. اگر برای مؤلفه‌ای $W_c>C$ باشد، نمونه ناموجه است؛ در غیر این صورت مسئله حاصل دقیقاً \lr{BPP} کلاسیک روی $|\mathcal{C}|$ قلم است. مؤلفه‌ها و وزن‌هایشان را می‌توان در زمان $O(|I|+|E|)$ محاسبه کرد.

\subsection{پیچیدگی و تقلیل}\label{subsec:lr_np_hardness}
نسخه تصمیم مسئله می‌پرسد آیا بسته‌بندی موجهی در حداکثر $K$ سطل وجود دارد. این مسئله در \lr{NP} است، زیرا با داشتن تخصیص اقلام می‌توان یکتایی تخصیص، ظرفیت‌ها، قیود یال‌ها و تعداد سطل‌ها را در زمان چندجمله‌ای بررسی کرد.

برای اثبات سختی قوی، یک نمونه دلخواه از نسخه تصمیم \lr{BPP} با اعداد صحیح، شامل $(I,(w_i)_{i\in I},C,K)$، را در نظر می‌گیریم. همان اقلام، وزن‌ها، ظرفیت و $K$ را نگه می‌داریم و فقط $E=\emptyset$ قرار می‌دهیم. نوشتن نمونه خروجی در زمانی خطی نسبت به طول ورودی انجام می‌شود و هیچ مقدار عددی تغییر یا بزرگ نمی‌شود. چون با $E=\emptyset$ هیچ قید چسبندگی وجود ندارد، پاسخ نمونه ساخته‌شده «بله» است اگر و تنها اگر پاسخ نمونه اصلی \lr{BPP} «بله» باشد. بنابراین این یک تقلیل چندجمله‌ایِ حفظ‌کننده اعداد است. از آنجا که \lr{BPP} از طریق \lr{3-Partition} قویاً \lr{NP-hard} است \cite{GareyJohnson1979}، مسئله حاضر نیز در بدترین حالت قویاً \lr{NP-hard} است؛ و با عضویت فوق در \lr{NP}، نسخه تصمیم آن قویاً \lr{NP-complete} است.

این نتیجه نباید چنین تفسیر شود که افزودن این نوع قید همواره مسئله را از \lr{BPP} سخت‌تر می‌کند. در واقع دو تبدیل چندجمله‌ای داریم: \lr{BPP} با انتخاب $E=\emptyset$ حالت خاص مسئله حاضر است، و هر نمونه مسئله حاضر با انقباض مؤلفه‌ها به یک \lr{BPP} کلاسیکِ هم‌ارز تبدیل می‌شود. بنابراین، برای همین تعریفِ صرفاً \lr{must-link}، تفاوت اصلی در پیش‌پردازش و اندازه اقلام تجمیع‌شده است، نه ایجاد یک کلاس بسته‌بندی ذاتاً سخت‌تر.

\subsection{مسائل مرتبط و کران‌های پایه}\label{subsec:32}
\begin{itemize}
    \item در \lr{Bin Packing with Conflicts}، یال‌ها هم‌مکانی را منع می‌کنند؛ ازاین‌رو contraction مؤلفه‌ها که در اینجا معتبر است، برای قیود تعارض معتبر نیست \cite{A1_Bacci2017,A2_DoronArad2020}.
    \item \lr{Cutting Stock} یک خانواده نزدیک از مسائل تخصیص تقاضا به مواد اولیه با ظرفیت محدود است، اما تقسیم‌پذیری اقلام، الگوهای برش و تابع هدف آن می‌تواند با مدل حاضر متفاوت باشد \cite{D1_Tanir2016}.
\end{itemize}

پس از contraction، دو کران پایین ابتدایی مستقیماً از \lr{BPP} کلاسیک به دست می‌آیند. اگر
\[
    L_{\mathrm{cap}}=\left\lceil\frac{\sum_{c\in\mathcal{C}}W_c}{C}\right\rceil
    \quad\text{و}\quad
    L_{>1/2}=\left|\{c\in\mathcal{C}:W_c>C/2\}\right|,
\]
آنگاه $\max\{L_{\mathrm{cap}},L_{>1/2}\}$ کران پایین معتبر است: اولی از مجموع ظرفیت لازم و دومی از این واقعیت ناشی می‌شود که دو اَبَرقلم با وزن بیش از $C/2$ در یک سطل جا نمی‌شوند. این کران‌ها لزوماً مقدار بهینه را تعیین نمی‌کنند و برابر دانستن آن‌ها با یک آزادسازی خطی قوی‌تر، بدون ارائه فرمول‌بندی و اثبات جداگانه، موجه نیست.
